\documentclass[11pt]{article}
\usepackage[margin=1in]{geometry}
\usepackage{booktabs}
\usepackage{longtable}
% public-pdf-table-layout-v1
\usepackage{array}
\usepackage{ragged2e}
\usepackage{xurl}
\usepackage{microtype}
\setlength{\emergencystretch}{4em}
\Urlmuskip=0mu plus 2mu\relax
\sloppy
\usepackage{hyperref}
% public-pdf-metadata-v1
\hypersetup{%
  pdftitle={Supplementary Document for the Regression CP Research Atlas},
  pdfauthor={Emre Tasar, Data Scientist},
  pdfsubject={Audited regression conformal prediction Research Atlas with evidence traceability, backend sensitivity analysis, and interpretation scope.},
  pdfkeywords={conformal prediction, regression, prediction intervals, Research Atlas, CQR, CV+, Venn-Abers, evidence graph},
  pdfcreator={LaTeX with hyperref},
  pdfproducer={pdfTeX},
  pdfencoding=auto,
  psdextra,
  pdflang={en-US}%
}
\usepackage[T1]{fontenc}
\usepackage{lmodern}
\title{Supplementary Document}
\author{Emre Tasar, Data Scientist}
\date{Working final-prose review draft; not final Research Document prose for public research report}
\begin{document}
\maketitle

Contact: detasar@gmail.com

\section{Supplement Reader Crosswalk}

This crosswalk explains how the supplement should be read with the main article. Each row identifies the main article surface, the supplement section that supplies the detailed audit evidence, the primary source artifact, and the stronger claim that remains beyond this study. The table is a navigation aid; it does not add experiments or promote a method.

\begingroup
\small
\setlength{\tabcolsep}{3pt}
\renewcommand{\arraystretch}{1.18}
\begin{longtable}{@{}>{\RaggedRight\arraybackslash}p{0.190\linewidth}>{\RaggedRight\arraybackslash}p{0.270\linewidth}>{\RaggedRight\arraybackslash}p{0.270\linewidth}>{\RaggedRight\arraybackslash}p{0.210\linewidth}@{}}
\toprule
Main article surface & Supplement support & Primary source artifact & Stronger claim requiring validation \\
\midrule
CQR/CV+ descriptive performance & S1 robustness diagnostics, S1b CQR backend sensitivity, and S2 post-selection diagnostics & \texttt{method\_\allowbreak{}selection\_\allowbreak{}robustness\_\allowbreak{}audit; cqr\_\allowbreak{}fixed\_\allowbreak{}vs\_\allowbreak{}model\_\allowbreak{}matched\_\allowbreak{}synthesis} & No final method selection, global best method, or deployment rule. \\
Venn-Abers bridge negative evidence & S2 bridge validation slice and S6 traceability & \texttt{dataset\_\allowbreak{}final\_\allowbreak{}gate\_\allowbreak{}post\_\allowbreak{}selection\_\allowbreak{}validation\_\allowbreak{}bridge\_\allowbreak{}results} & No validated Venn-Abers regression interval claim. \\
Bounded-support validity & S3 bounded-support endpoint policy & \texttt{bounded\_\allowbreak{}support\_\allowbreak{}endpoint\_\allowbreak{}closure\_\allowbreak{}audit} & No endpoint or bounded-support validity claim. \\
Population-level group inference & S4 group diagnostics & \texttt{group inference\_\allowbreak{}population\_\allowbreak{}readiness\_\allowbreak{}audit} & No population-level group inference claim. \\
Duplicate and leakage integrity & S5 duplicate and integrity caveats & \texttt{cross\_\allowbreak{}run\_\allowbreak{}integrity\_\allowbreak{}audit} & No claim strengthening by hiding duplicate or split-integrity caveats. \\
Knowledge-graph traceability & S6 traceability & \texttt{knowledge\_\allowbreak{}graph\_\allowbreak{}quality\_\allowbreak{}summary} & The KG and site support navigation and source tracing. \\
\bottomrule
\end{longtable}
\endgroup

\section{Supplement Reading Protocol}

The supplement is intentionally broader than the main article. Each section adds audit detail, but the extra detail must be read through a specific boundary. The protocol below states what a reviewer should check, which evidence metric anchors the check, the allowed use of that evidence, and the overclaim that remains beyond this study.

\begingroup
\small
\setlength{\tabcolsep}{3pt}
\renewcommand{\arraystretch}{1.18}
\begin{longtable}{@{}>{\RaggedRight\arraybackslash}p{0.150\linewidth}>{\RaggedRight\arraybackslash}p{0.220\linewidth}>{\RaggedRight\arraybackslash}p{0.220\linewidth}>{\RaggedRight\arraybackslash}p{0.200\linewidth}>{\RaggedRight\arraybackslash}p{0.150\linewidth}@{}}
\toprule
Section & Reviewer check & Evidence metric & Allowed use & Beyond this study use \\
\midrule
S1 & Check whether the method-selection signal is robust enough to describe, but not promote. & CQR common-cell wins 58 of 94 cells; bootstrap selection count is cqr=1,000. & Within this retrospective and imbalanced experiment surface, the fixed-GBM CQR pipeline was selected most often under the coverage-gated interval-score rule; Mondrian calibration and CV+ were secondary practical candidates. & CQR final-selection and universal-deployment readings require separate validation. \\
S1b & Check whether CQR's signal survives a backend-confound diagnostic. & Model-matched CQR completed 4,564 rows; paired cells=224; selected cells are fixed\_\allowbreak{}gbm\_\allowbreak{}cqr=116, model\_\allowbreak{}matched\_\allowbreak{}cqr=71, neither=37. & Use as a backend sensitivity check for the experiment-scoped CQR signal. & The check remains backend-sensitivity evidence for the CQR pipeline signal. \\
S2 & Check whether post-selection and bridge diagnostics preserve the same boundary. & Post-selection diagnostic counts are cqr=18 and mondrian\_\allowbreak{}abs=7; the bridge slice has 45 completed atomic runs. & Use the slice as diagnostic stress-test evidence. & The slice remains diagnostic evidence rather than final selection or validated Venn-Abers interval evidence. \\
S3 & Check whether bounded-support endpoint behavior authorizes a validity claim. & 0 bounded-support-validity-ready bundles and 11 raw endpoint-excursion bundles. & Report endpoint policy closure and retained domain caveats. & Bounded-support and endpoint-validity claims require separate validation. \\
S4 & Check whether group diagnostics define a population-group-inference estimand. & 187 pairwise group comparisons, 0 population-group-inference-ready bundles, and 0 population-estimand-declared bundles. & Use group rows as heterogeneity diagnostics. & Population-level and deployment group-inference claims require separate validation. \\
S5 & Check whether duplicate and leakage controls strengthen or limit interpretation. & 29 duplicate actions, 46 quarantined actions, 0 unquarantined actions, and hard leakage not detected in scanned artifacts. & Retain integrity caveats as part of the scientific record. & Integrity caveats remain part of the scientific record. \\
S6 & Check whether traceability artifacts are review infrastructure or public outputs. & 3,643 KG nodes, 21,019 edges, 0 isolated nodes, and evidence scope recorded in the manifest. & Use the KG and package for working claim tracing. & Use the KG and site to inspect how claims connect to evidence. \\
\bottomrule
\end{longtable}
\endgroup

\section{S1. Method Selection Robustness Diagnostics}

This section supports the main article by separating descriptive method behavior from final method selection. In conformal prediction, \texttt{1 - alpha} is the target coverage level and \texttt{alpha} is the target miscoverage rate \cite{lei2017distribution_free_regression}. CQR builds an interval from lower and upper quantile models before calibration \cite{romano2019conformalized_quantile_regression}, while CV+ and jackknife-style methods use out-of-fold or leave-one-out predictions \cite{barber2020jackknife_plus,kim2020jackknife_after_bootstrap}.

\begingroup
\small
\setlength{\tabcolsep}{3pt}
\renewcommand{\arraystretch}{1.18}
\begin{longtable}{@{}>{\RaggedRight\arraybackslash}p{0.250\linewidth}>{\RaggedRight\arraybackslash}p{0.340\linewidth}>{\RaggedRight\arraybackslash}p{0.350\linewidth}@{}}
\toprule
Diagnostic item & Value & Boundary \\
\midrule
Candidate methods & 3 & CQR, CV+, and Mondrian absolute-residual calibration only \\
Common dataset-alpha cells & 94 & Shared comparison cells \\
Common-cell selected diagnostic method & \texttt{cqr} & Not a final method selection \\
Common-cell counts & cqr=58; cv\_\allowbreak{}plus=15; mondrian\_\allowbreak{}abs=21 & Descriptive counts only \\
Bootstrap replicates & 1,000 & Selection stability diagnostic \\
Bootstrap counts & cqr=1,000 & Experiment-scoped interpretation \\
Bootstrap primary selection rate & 1.0000 & Stability estimate, not final selection \\
Bootstrap primary selection-rate diagnostic band & 0.9962 to 1.0000 & Precision check for the diagnostic rate \\
Leave-one-dataset retention & 1.0000 & Diagnostic robustness \\
Leave-one-alpha retention & 1.0000 & Diagnostic robustness \\
Final selection claim & \texttt{beyond this study} & Final selection remains beyond this study \\
Pairwise shared cells & 94 & Inferential support, not a claim upgrade \\
Main-result primary win rate & 0.6778 & Descriptive rate with final claim beyond this study \\
Main-result primary win-rate diagnostic band & 0.5757 to 0.7653 & Uncertainty around the descriptive rate \\
Robustness common-cell primary win rate & 0.6170 & Common-cell diagnostic only \\
Robustness common-cell diagnostic band & 0.5160 to 0.7089 & Precision check, not claim scope \\
\bottomrule
\end{longtable}
\endgroup

Reader note: this section explains why the fixed-GBM CQR pipeline, Mondrian calibration, and CV+ can be reported as practical candidates within this experiment while the final-selection gate remains beyond this study. The diagnostic bands summarize uncertainty around the observed rates; they do not change the claim scope of the final-selection claim.

\subsection{S1b. CQR Backend Sensitivity Check}

The completed model-matched CQR rerun checks whether the historical fixed-GBM CQR backend explains the CQR signal. It preserves the experiment-scoped interpretation and does not create a method-selection claim.

\begingroup
\small
\setlength{\tabcolsep}{3pt}
\renewcommand{\arraystretch}{1.18}
\begin{longtable}{@{}>{\RaggedRight\arraybackslash}p{0.250\linewidth}>{\RaggedRight\arraybackslash}p{0.340\linewidth}>{\RaggedRight\arraybackslash}p{0.350\linewidth}@{}}
\toprule
Backend diagnostic item & Value & Boundary \\
\midrule
Fixed-GBM CQR completed rows & 4,564 & Historical comparator rows \\
Model-matched CQR completed rows & 4,564 & Completed rerun rows \\
Paired dataset-alpha-model-family cells & 224 & Direct CQR backend comparison cells \\
Fixed-GBM CQR selected cells & 116 & Coverage-eligible interval-score selections \\
Model-matched CQR selected cells & 71 & Coverage-eligible interval-score selections \\
Neither coverage-eligible variant & 37 & Cells where both CQR variants fail the coverage-eligibility rule \\
Method-selection claim supported & \texttt{False} & No deployment guidance follows from this table \\
\bottomrule
\end{longtable}
\endgroup

\section{S2. Post-Selection Validation Diagnostics}

Post-selection validation asks whether the descriptive pattern persists in a held-out validation-style slice after the candidate methods have already been identified. The supplement treats this as a diagnostic stress test. It does not convert the observed counts into a final method selection.

\begingroup
\small
\setlength{\tabcolsep}{3pt}
\renewcommand{\arraystretch}{1.18}
\begin{longtable}{@{}>{\RaggedRight\arraybackslash}p{0.250\linewidth}>{\RaggedRight\arraybackslash}p{0.340\linewidth}>{\RaggedRight\arraybackslash}p{0.350\linewidth}@{}}
\toprule
Diagnostic item & Value & Boundary \\
\midrule
Validation datasets & 5 & Narrow validation scope \\
Validation dataset-alpha cells & 25 & Shared cells only \\
Completed validation atomic runs & 225 & Completed diagnostic rows \\
Validation diagnostic counts & cqr=18; mondrian\_\allowbreak{}abs=7 & No final method selection \\
Validation primary win rate & 0.7200 & Post-selection diagnostic rate \\
Validation primary win-rate diagnostic band & 0.5242 to 0.8572 & Wide uncertainty remains visible \\
Feature leakage violations & 0 & Leakage sidecars clean in this slice \\
Width pathology rows & 6 & Retained caveat evidence \\
Bridge validation datasets & 1 & Dataset-final-gate bridge slice \\
Bridge validation atomic runs & 45 & Auxiliary diagnostic evidence \\
Bridge validation diagnostic counts & cqr=3; cv\_\allowbreak{}plus=1; mondrian\_\allowbreak{}abs=1 & Bridge-slice diagnostics only \\
Venn-Abers bridge undercoverage runs & 14 & Negative bridge evidence \\
\bottomrule
\end{longtable}
\endgroup

Reader note: post-selection diagnostics are stress tests of a descriptive pattern. They do not license language such as final method selection, globally best conformal method, or validated Venn-Abers regression interval. The Venn-Abers rows are especially narrow: they concern the evaluated bridge, while predictive-distribution and generalized Venn-Abers work remain separate literature objects \cite{nouretdinov2018ivapd,nouretdinov2024ivapd_applications,vanderlaan2025generalized_venn_abers,petej2026inductive_venn_abers_regressors}.

\section{S3. Bounded-Support Endpoint Policy}

Bounded-support checks ask whether prediction interval endpoints respect the natural domain of the target. This is different from ordinary marginal coverage. Weighted conformal ideas can address some covariate-shift settings \cite{tibshirani2020covariate_shift}, but the present endpoint policy is an empirical domain-validity gate. The current gate blocks bounded-support validity claims.

\begingroup
\small
\setlength{\tabcolsep}{3pt}
\renewcommand{\arraystretch}{1.18}
\begin{longtable}{@{}>{\RaggedRight\arraybackslash}p{0.250\linewidth}>{\RaggedRight\arraybackslash}p{0.340\linewidth}>{\RaggedRight\arraybackslash}p{0.350\linewidth}@{}}
\toprule
Diagnostic item & Value & Boundary \\
\midrule
Bounded-support bundles & 15 & Research Document-candidate bundles \\
Bounded-support datasets & 6 & Current endpoint-audit scope \\
Raw endpoint-excursion bundles & 11 & Domain blocker retained \\
Post-handling validated bundles & 15 & Policy triage complete \\
Clean or not-applicable bundles & 4 & Not enough for global validity \\
Validity-ready bundles & 0 & No bounded-support validity claim \\
Positive-claim-ready bundles & 0 & Positive endpoint claim beyond this study \\
Endpoint support status counts & beyond this study\_\allowbreak{}natural\_\allowbreak{}domain\_\allowbreak{}endpoint\_\allowbreak{}excursions=11; clean\_\allowbreak{}no\_\allowbreak{}natural\_\allowbreak{}domain\_\allowbreak{}endpoint\_\allowbreak{}excursions=3; not\_\allowbreak{}applicable\_\allowbreak{}unbounded\_\allowbreak{}target\_\allowbreak{}endpoint\_\allowbreak{}hygiene\_\allowbreak{}recorded=1 & Support-status inventory, not a validity upgrade \\
\bottomrule
\end{longtable}
\endgroup

Reader note: bounded-support evidence is stricter than ordinary coverage evidence. A method can have acceptable empirical coverage and still fail a natural-domain endpoint gate. The supplement therefore reports endpoint closure as a limitation, not as a post-processing success story.

\section{S4. Group Diagnostics}

Group coverage diagnostics compare interval behavior across observed groups. They are useful for auditing heterogeneity, but they are not population-level group inference claims unless the population estimand, protected-attribute scope, sampling design, and multiplicity treatment support that claim.

\begingroup
\small
\setlength{\tabcolsep}{3pt}
\renewcommand{\arraystretch}{1.18}
\begin{longtable}{@{}>{\RaggedRight\arraybackslash}p{0.250\linewidth}>{\RaggedRight\arraybackslash}p{0.340\linewidth}>{\RaggedRight\arraybackslash}p{0.350\linewidth}@{}}
\toprule
Diagnostic item & Value & Boundary \\
\midrule
Group diagnostic bundles & 15 & Diagnostic group scope \\
Group diagnostic datasets & 6 & Current diagnostic scope \\
Bootstrap replicates & 500 & Gap-uncertainty diagnostic \\
Group counts recorded & 15 & Required diagnostic metadata \\
Coverage by group recorded & 15 & Group comparison evidence \\
Width by group recorded & 15 & Interval-size comparison evidence \\
Gap uncertainty recorded & 15 & Multiplicity-aware caveat evidence \\
Pairwise group comparisons & 187 & Multiplicity scope, not group-inference proof \\
Population-estimand-declared bundles & 0 & Required for population-level group inference, absent here \\
Protected-attribute-scope-declared bundles & 0 & Required scope declaration, absent here \\
Weighted-estimand-applied bundles & 0 & No weighted population estimand applied \\
Sampling-weight-policy-declared bundles & 15 & Policy exists but does not open group-inference claim \\
Population-group-inference-ready bundles & 0 & No population-level group inference claim \\
\bottomrule
\end{longtable}
\endgroup

Reader note: group rows can show where interval behavior differs, but group inference is an estimand-level claim. Because the current population estimand and protected-attribute scope are not declared at the bundle level, the correct supplement reading is diagnostic heterogeneity, not population-level group inference.

\section{S5. Duplicate And Integrity Caveats}

Duplicate and split-integrity checks protect the interpretation of model comparisons. If repeated rows, row signatures, or model-visible duplicates can affect train/calibration/test separation, the result must retain the caveat rather than become stronger prose.

\begingroup
\small
\setlength{\tabcolsep}{3pt}
\renewcommand{\arraystretch}{1.18}
\begin{longtable}{@{}>{\RaggedRight\arraybackslash}p{0.250\linewidth}>{\RaggedRight\arraybackslash}p{0.340\linewidth}>{\RaggedRight\arraybackslash}p{0.350\linewidth}@{}}
\toprule
Diagnostic item & Value & Boundary \\
\midrule
Duplicate actions & 29 & Caveat inventory scope \\
Open duplicate actions & 0 & Closure evidence only \\
Paired duplicate comparison rows & 1,043 & Sensitivity evidence \\
Quarantined actions & 46 & Main-claim promotion beyond this study \\
Unquarantined actions & 0 & No untracked duplicate action \\
Cross-run leakage status & \texttt{hard\_\allowbreak{}leakage\_\allowbreak{}not\_\allowbreak{}detected\_\allowbreak{}in\_\allowbreak{}scanned\_\allowbreak{}artifacts} & Hard leakage not detected in scanned artifacts \\
Unsupported claim hits & 0 & Claim scan clean for this audit \\
Cross-run risk counts & medium=39; pass=109 & Risk inventory retained \\
Cross-run caveat counts & duplicate\_\allowbreak{}cluster\_\allowbreak{}plus\_\allowbreak{}family\_\allowbreak{}internal\_\allowbreak{}fold\_\allowbreak{}caveat=17; duplicate\_\allowbreak{}signature\_\allowbreak{}cross\_\allowbreak{}split\_\allowbreak{}caveat=14; model\_\allowbreak{}visible\_\allowbreak{}signature\_\allowbreak{}cross\_\allowbreak{}split\_\allowbreak{}caveat=15 & Caveats constrain interpretation \\
\bottomrule
\end{longtable}
\endgroup

Reader note: the clean unsupported-claim scan does not erase duplicate or split caveats. It means the scanned artifacts did not contain unsupported claim language; the caveats still travel with the empirical interpretation.

\section{S6. Traceability And Release State}

The supplement is linked to a knowledge graph snapshot with 3,643 nodes, 21,019 edges, and 0 isolated nodes. The graph connects report sections, source artifacts, citations, interpretation scope, and quality checks. It remains an internal evidence artifact until the sterile repository and release review are complete.

Venn-Abers evidence is retained as part of the negative and boundary evidence for this study \cite{nouretdinov2018ivapd}. The supplement does not claim that the current interval bridge validates Venn-Abers regression intervals.

\section{Interpretation Scope}

\begin{itemize}
\item This is a working final-prose supplementary review draft, not final Research Document prose for public research report.
\item Method evidence is diagnostic; no final method selection or deployment rule is established.
\item Bounded-support endpoint results block validity claims.
\item Group diagnostics do not establish population-level group inference.
\item Duplicate and integrity caveats are retained rather than hidden.
\end{itemize}

\section{References}

\begin{itemize}
\item \texttt{@barber2020jackknife\_\allowbreak{}plus}: https://arxiv.org/abs/1905.02928
\item \texttt{@kim2020jackknife\_\allowbreak{}after\_\allowbreak{}bootstrap}: https://arxiv.org/abs/2002.09025
\item \texttt{@lei2017distribution\_\allowbreak{}free\_\allowbreak{}regression}: https://arxiv.org/abs/1604.04173
\item \texttt{@nouretdinov2018ivapd}: https://proceedings.mlr.press/v91/nouretdinov18a.html
\item \texttt{@nouretdinov2024ivapd\_\allowbreak{}applications}: https://proceedings.mlr.press/v230/nouretdinov24a.html
\item \texttt{@petej2026inductive\_\allowbreak{}venn\_\allowbreak{}abers\_\allowbreak{}regressors}: https://arxiv.org/html/2605.06646v1
\item \texttt{@romano2019conformalized\_\allowbreak{}quantile\_\allowbreak{}regression}: https://arxiv.org/abs/1905.03222
\item \texttt{@tibshirani2020covariate\_\allowbreak{}shift}: https://arxiv.org/abs/1904.06019
\item \texttt{@vanderlaan2025generalized\_\allowbreak{}venn\_\allowbreak{}abers}: https://proceedings.mlr.press/v267/van-der-laan25a.html
\end{itemize}

\section{Source Artifacts}

\begin{itemize}
\item \texttt{main\_\allowbreak{}article\_\allowbreak{}draft}: \texttt{experiments/\allowbreak{}regression/\allowbreak{}Research Document/\allowbreak{}main\_\allowbreak{}article\_\allowbreak{}draft.\allowbreak{}json}
\item \texttt{individual\_\allowbreak{}experiment\_\allowbreak{}report\_\allowbreak{}draft}: \texttt{experiments/\allowbreak{}regression/\allowbreak{}Research Document/\allowbreak{}individual\_\allowbreak{}experiment\_\allowbreak{}report\_\allowbreak{}draft.\allowbreak{}json}
\item \texttt{article\_\allowbreak{}supplement\_\allowbreak{}blueprint\_\allowbreak{}alignment}: \texttt{experiments/\allowbreak{}regression/\allowbreak{}Research Document/\allowbreak{}article\_\allowbreak{}supplement\_\allowbreak{}blueprint\_\allowbreak{}alignment.\allowbreak{}json}
\item \texttt{method\_\allowbreak{}selection\_\allowbreak{}robustness\_\allowbreak{}audit}: \texttt{experiments/\allowbreak{}regression/\allowbreak{}reports/\allowbreak{}methodology\_\allowbreak{}sanity\_\allowbreak{}audit\_\allowbreak{}20260627/\allowbreak{}method\_\allowbreak{}selection\_\allowbreak{}robustness\_\allowbreak{}audit.\allowbreak{}json}
\item \texttt{method\_\allowbreak{}selection\_\allowbreak{}inferential\_\allowbreak{}audit}: \texttt{experiments/\allowbreak{}regression/\allowbreak{}reports/\allowbreak{}methodology\_\allowbreak{}sanity\_\allowbreak{}audit\_\allowbreak{}20260627/\allowbreak{}method\_\allowbreak{}selection\_\allowbreak{}inferential\_\allowbreak{}audit.\allowbreak{}json}
\item \texttt{method\_\allowbreak{}selection\_\allowbreak{}post\_\allowbreak{}selection\_\allowbreak{}validation\_\allowbreak{}results}: \texttt{experiments/\allowbreak{}regression/\allowbreak{}reports/\allowbreak{}methodology\_\allowbreak{}sanity\_\allowbreak{}audit\_\allowbreak{}20260627/\allowbreak{}method\_\allowbreak{}selection\_\allowbreak{}post\_\allowbreak{}selection\_\allowbreak{}validation\_\allowbreak{}results.\allowbreak{}json}
\item \texttt{dataset\_\allowbreak{}final\_\allowbreak{}gate\_\allowbreak{}post\_\allowbreak{}selection\_\allowbreak{}validation\_\allowbreak{}bridge\_\allowbreak{}results}: \texttt{experiments/\allowbreak{}regression/\allowbreak{}reports/\allowbreak{}methodology\_\allowbreak{}sanity\_\allowbreak{}audit\_\allowbreak{}20260627/\allowbreak{}dataset\_\allowbreak{}final\_\allowbreak{}gate\_\allowbreak{}post\_\allowbreak{}selection\_\allowbreak{}validation\_\allowbreak{}bridge\_\allowbreak{}results.\allowbreak{}json}
\item \texttt{bounded\_\allowbreak{}support\_\allowbreak{}endpoint\_\allowbreak{}closure\_\allowbreak{}audit}: \texttt{experiments/\allowbreak{}regression/\allowbreak{}reports/\allowbreak{}methodology\_\allowbreak{}sanity\_\allowbreak{}audit\_\allowbreak{}20260627/\allowbreak{}bounded\_\allowbreak{}support\_\allowbreak{}endpoint\_\allowbreak{}closure\_\allowbreak{}audit.\allowbreak{}json}
\item \texttt{bounded\_\allowbreak{}support\_\allowbreak{}positive\_\allowbreak{}validation\_\allowbreak{}protocol}: \texttt{experiments/\allowbreak{}regression/\allowbreak{}Research Document/\allowbreak{}bounded\_\allowbreak{}support\_\allowbreak{}positive\_\allowbreak{}validation\_\allowbreak{}protocol.\allowbreak{}json}
\item \texttt{group inference\_\allowbreak{}group\_\allowbreak{}diagnostic\_\allowbreak{}audit}: \texttt{experiments/\allowbreak{}regression/\allowbreak{}reports/\allowbreak{}methodology\_\allowbreak{}sanity\_\allowbreak{}audit\_\allowbreak{}20260627/\allowbreak{}group inference\_\allowbreak{}group\_\allowbreak{}diagnostic\_\allowbreak{}audit.\allowbreak{}json}
\item \texttt{group inference\_\allowbreak{}population\_\allowbreak{}readiness\_\allowbreak{}audit}: \texttt{experiments/\allowbreak{}regression/\allowbreak{}reports/\allowbreak{}methodology\_\allowbreak{}sanity\_\allowbreak{}audit\_\allowbreak{}20260627/\allowbreak{}group inference\_\allowbreak{}population\_\allowbreak{}readiness\_\allowbreak{}audit.\allowbreak{}json}
\item \texttt{group inference\_\allowbreak{}group\_\allowbreak{}multiplicity\_\allowbreak{}scope}: \texttt{experiments/\allowbreak{}regression/\allowbreak{}Research Document/\allowbreak{}group inference\_\allowbreak{}group\_\allowbreak{}multiplicity\_\allowbreak{}scope.\allowbreak{}json}
\item \texttt{duplicate\_\allowbreak{}sensitivity\_\allowbreak{}closure\_\allowbreak{}audit}: \texttt{experiments/\allowbreak{}regression/\allowbreak{}reports/\allowbreak{}methodology\_\allowbreak{}sanity\_\allowbreak{}audit\_\allowbreak{}20260627/\allowbreak{}duplicate\_\allowbreak{}sensitivity\_\allowbreak{}closure\_\allowbreak{}audit.\allowbreak{}json}
\item \texttt{duplicate\_\allowbreak{}content\_\allowbreak{}quarantine\_\allowbreak{}audit}: \texttt{experiments/\allowbreak{}regression/\allowbreak{}reports/\allowbreak{}methodology\_\allowbreak{}sanity\_\allowbreak{}audit\_\allowbreak{}20260627/\allowbreak{}duplicate\_\allowbreak{}content\_\allowbreak{}quarantine\_\allowbreak{}audit.\allowbreak{}json}
\item \texttt{cross\_\allowbreak{}run\_\allowbreak{}integrity\_\allowbreak{}audit}: \texttt{experiments/\allowbreak{}regression/\allowbreak{}reports/\allowbreak{}methodology\_\allowbreak{}sanity\_\allowbreak{}audit\_\allowbreak{}20260627/\allowbreak{}cross\_\allowbreak{}run\_\allowbreak{}integrity\_\allowbreak{}audit.\allowbreak{}json}
\item \texttt{publication\_\allowbreak{}citation\_\allowbreak{}registry}: \texttt{experiments/\allowbreak{}regression/\allowbreak{}Research Document/\allowbreak{}publication\_\allowbreak{}citation\_\allowbreak{}registry.\allowbreak{}json}
\item \texttt{knowledge\_\allowbreak{}graph\_\allowbreak{}quality\_\allowbreak{}summary}: \texttt{experiments/\allowbreak{}regression/\allowbreak{}reports/\allowbreak{}knowledge\_\allowbreak{}graph\_\allowbreak{}quality/\allowbreak{}quality\_\allowbreak{}summary.\allowbreak{}json}
\item \texttt{cqr\_\allowbreak{}fixed\_\allowbreak{}vs\_\allowbreak{}model\_\allowbreak{}matched\_\allowbreak{}synthesis}: \texttt{experiments/\allowbreak{}regression/\allowbreak{}reports/\allowbreak{}model\_\allowbreak{}matched\_\allowbreak{}cqr\_\allowbreak{}rerun\_\allowbreak{}plan/\allowbreak{}cqr\_\allowbreak{}fixed\_\allowbreak{}vs\_\allowbreak{}model\_\allowbreak{}matched\_\allowbreak{}synthesis.\allowbreak{}json}
\end{itemize}

\bibliographystyle{plain}
\bibliography{references}

\end{document}
